\begin{proofbox}
	\textbf{\textcolor{CProof}{证明：严格递增的充要条件}}

	记 $L=\lim_{x\to x_0^+}h(x)\in\mathbb{R}$。

	\textbf{必要性。} 若 $f$ 在 $\mathbb{R}$ 上严格递增，将定义限制到两段，各自必严格递增。
	又对每个 $x>x_0$，有 $g(x_0)=f(x_0)<f(x)=h(x)$。令 $x\to x_0^+$，得到 $g(x_0)\le L$。
	注意取极限后只能保证\emph{不大于}，不能保证严格小于。

	\medskip
	\textbf{充分性。} 假设 $g,h$ 在各自区间严格递增，且 $g(x_0)\le L$。
	任取 $x_1<x_2$，分三种情况：
	\begin{enumerate}[leftmargin=*, itemsep=0.35em]
		\item 若 $x_1<x_2\le x_0$，由 $g$ 严格递增，得 $f(x_1)<f(x_2)$。
		\item 若 $x_0<x_1<x_2$，由 $h$ 严格递增，得 $f(x_1)<f(x_2)$。
		\item 若 $x_1\le x_0<x_2$，则 $g(x_1)\le g(x_0)\le L$。取 $t\in(x_0,x_2)$，由 $h$ 严格递增且 $L\le h(t)<h(x_2)$，得 $L<h(x_2)$。因此
		      \[
			      f(x_1)=g(x_1)\le g(x_0)\le L<h(x_2)=f(x_2).
		      \]
	\end{enumerate}
	三种情况均有 $f(x_1)<f(x_2)$，故 $f$ 在 $\mathbb{R}$ 上严格递增。

	\medskip
	\textbf{\textcolor{CProof}{严格递减：}}
	同理，将段内方向与全部不等号反向；跨段时利用 $h(x_2)<L$，得到
	$f(x_1)=g(x_1)\ge g(x_0)\ge L>h(x_2)=f(x_2)$。
\end{proofbox}